\section{Introducción}

\subsection{Contexto y Motivación}
El presente trabajo describe el desarrollo de \textbf{Pokemonad}, un motor de juego de combate por turnos inspirado en la mecánica clásica de Pokémon, construido íntegramente utilizando el lenguaje de programación funcional Haskell. 

La motivación principal del proyecto no radica en la réplica exhaustiva de las complejas reglas del juego original, sino en utilizar este dominio interactivo como vehículo para explorar y resolver problemas de ingeniería de software dentro del paradigma funcional. Tras la etapa de ideación y evaluación de alcance, el enfoque del proyecto se centró en dos grandes desafíos técnicos: la implementación de una arquitectura de red descentralizada Peer-to-Peer (P2P) concurrente y segura, y el desarrollo de una Inteligencia Artificial no trivial para el modo de un solo jugador. 

Para lograr la profundidad requerida en estos subsistemas, se tomó la decisión arquitectónica de mantener la lógica de dominio del combate en su forma más pura y básica. Las interacciones se restringieron exclusivamente a acciones de ataque directo (\textit{FIGHT}) y cambio de criaturas (\textit{SWITCH}), omitiendo alteraciones de estado, modificadores de estadísticas y habilidades pasivas. Esta reducción deliberada del alcance permite concentrar el esfuerzo de desarrollo en la gestión del estado concurrente y en los árboles de decisión del bot oponente.

\subsection{Objetivos del Proyecto}
El objetivo general del trabajo es demostrar la aplicación práctica de conceptos avanzados de programación funcional, tales como la inmutabilidad, la evaluación de funciones puras, y el manejo de efectos a través de mónadas, en la construcción de un sistema de software interactivo y distribuido. 

Los objetivos específicos, acordados y delimitados para el alcance de este trabajo, incluyen:
\begin{itemize}
    \item \textbf{Desarrollo de una biblioteca P2P independiente:} Diseñar un módulo de comunicación de bajo nivel mediante sockets TCP puro. Este subsistema gestiona la recepción asincrónica de flujos de bytes (con serialización binaria en tramas) y delega la sincronización de estados concurrentes mediante Memoria Transaccional en Software (STM, usando \texttt{TQueue}). Este módulo opera de forma completamente desacoplada de las reglas de \textit{Pokemonad}, constituyendo una biblioteca genérica.
    \item \textbf{Implementación de un oponente controlado por IA:} Construir un agente inteligente para el modo de un jugador que supere la toma de decisiones heurística trivial. Se implementó un modelo basado en \textbf{Aprendizaje por Refuerzo (Reinforcement Learning)}, utilizando una aproximación lineal para estimar los valores Q (Q-Values) a partir de la extracción de características del entorno (ventaja de tipos, proyecciones de daño y viabilidad de cambio de criaturas). La selección de acciones se rige por una política $\epsilon$-greedy parametrizada según la dificultad del oponente.
    \item \textbf{Integración de una Interfaz Gráfica (GUI):} Proveer una capa visual interactiva utilizando la biblioteca \texttt{gloss}. Esta herramienta permite renderizar gráficos 2D de forma declarativa, exponiendo el estado inmutable del motor en tiempo real e integrándose fluidamente con el bucle principal de la aplicación.
\end{itemize}

\subsection{Integrantes y Metodología de Trabajo}
El proyecto fue diseñado y desarrollado íntegramente por Francisco Quian Blanco y Theo Stanfield. 

Dada la complejidad inherente de orquestar un hilo de red asincrónico junto con un hilo de interfaz gráfica (UI) evitando condiciones de carrera, se adoptó una metodología de desarrollo colaborativa y altamente comunicativa. La división de tareas se gestionó de manera dinámica apoyándose en un sistema de control de versiones. 

Para los componentes críticos de la arquitectura —particularmente el aislamiento de la biblioteca P2P y la implementación de la IA— se priorizó el trabajo conjunto mediante sesiones de \textit{pair programming}. Para el desarrollo de módulos individuales, se establecieron revisiones de código cruzadas (\textit{code reviews}). Esta metodología garantiza no solo el avance equilibrado del proyecto, sino que asegura que ambos integrantes mantengan un conocimiento total y un manejo profundo de cada aspecto del código fuente producido, cumpliendo rigurosamente con los requisitos de evaluación.

\clearpage